\bookchapter{פרק 67: ליטוש והוצאה לאור — בדיקות, ניפוי באגים ועלייה ל-App Store}

\chapterintro{האפליקציה נבנתה — ומה עכשיו? הפרק הזה עונה על שלושה עניינים בדרך מ“זה רץ” אל “זה מוכן למסירה”: איך לנפות באגים ביעילות, איך לכתוב בדיקות שמונעות נסיגות, ואיך להעלות את האפליקציה ל-App Store כדי שכל העולם יוכל להוריד. העלייה לחנות היא לא רק טקס — זהו קו הפרדה בין יצירה למוצר.}

\booksection{שלושת כלי הניפוי: print, נקודות עצירה ותצוגות מקדימות}

\begin{center}
\begin{tabularx}{\linewidth}{|X|X|X|}
\hline
\textbf{כלי} & \textbf{מה הוא עושה} & \textbf{איך משתמשים} \tabularnewline
\hline
print + הקונסולה & צפייה בלוגי הריצה & בקוד כותבים print("ערך=\textbackslash{}(x)") ורואים את הפלט בקונסולה שבתחתית Xcode \tabularnewline
ניפוי עם נקודות עצירה & עוצרים את התוכנית בשורה מסוימת ורואים ערכי משתנים & לוחצים פעם אחת ליד מספר השורה ומופיע סימן כחול; הריצה נעצרת שם ואפשר להתקדם צעד-צעד \tabularnewline
תצוגה מקדימה של SwiftUI + מנפה השגיאות & צפייה בזמן אמת בממשק ובהיררכיית התצוגות & בד ה-\#Preview + בודק ה-View Hierarchy בזמן ריצה \tabularnewline
\hline
\end{tabularx}
\end{center}

\begin{center}\small\color{gray}שלושת כלי הניפוי\end{center}

\begin{tcbinfo}{טיפ: איך קוראים קריסה?}
כש-Swift קורס, הקונסולה מדפיסה “Fatal error” ואת הסיבה, ו-Xcode עוצר בשורה הבעייתית ומדגיש אותה. היסטוריית קריסות ממכשיר אמיתי אפשר לראות ב-Xcode Window Organizer Crashes (נדרש חשבון בתשלום). בחיפוש, תעדפו את מילות המפתח “המשפט הראשון של סיבת השגיאה + swiftui” — רק שאלה מדויקת מניבה תשובה מדויקת.

\end{tcbinfo}

\booksection{בדיקות: ביטוח לקוד}

בדיקה ידנית לא אמינה — מתקנים פיצ׳ר אחד ואחר שובר בשקט. בדיקות הן גם “תרופת החרטה” וגם “שריון חץ”. ב-iOS משתמשים בפריימוורק XCTest, עם שני סוגים:

\begin{center}
\begin{tabularx}{\linewidth}{|X|X|X|X|}
\hline
\textbf{רמה} & \textbf{מה נבדק} & \textbf{Target} & \textbf{מהירות} \tabularnewline
\hline
בדיקות יחידה & לוגיקה טהורה (חישובים, בדיקות תקינות) & MyFirstAppTests & שניות \tabularnewline
בדיקות UI & התנהגות ממשק אמיתית & MyFirstAppUITests & דקות \tabularnewline
בדיקה ידנית & תחושות וחוויה של אדם אמיתי & —— & תמיד נדרשת \tabularnewline
\hline
\end{tabularx}
\end{center}

\begin{center}\small\color{gray}רמות הבדיקה ב-iOS (בפרויקט התבנית כבר מוגדרים שני Test Target)\end{center}

\begin{lstlisting}[language=Swift]
import XCTest

// לחלץ את הלוגיקה של "המרת טקסט הסכום למספר" לפונקציה טהורה
func parseAmount(_ input: String) -> Double? {
    guard let v = Double(input.trimmingCharacters(in: .whitespaces)),
          v > 0 else { return nil }
    return v
}

final class AmountTests: XCTestCase {
    func testEmptyInputReturnsNil() {
        XCTAssertNil(parseAmount(""))
    }
    func testNegativeReturnsNil() {
        XCTAssertNil(parseAmount("-5"))
    }
    func testNormalValueParses() {
        XCTAssertEqual(parseAmount("12.5"), 12.5)
    }
}
\end{lstlisting}

הרצת בדיקות: CmdU. רק אחרי שהכל ירוק שווה לשנות קוד — אז יש לכם רגליים על הקרקע. את בדיקות התקינות של הסכומים ואת חישובי הסיכום בספר החשבונות כדאי להצמיד לבדיקות יחידה.

\booksection{חתימה ואריזה: נותנים לאפליקציה “חותמת”}

כל אפליקציית iOS חייבת להיות חתומה כדי להתקין — זה מוכיח ש“החבילה הזאת יצאה ממך ולא זויפה”. היררכיית ההסתמכות:

\begin{itemize}
  \item שלב הלמידה: Apple ID חינמי + חתימה אוטומטית של Xcode (Personal Team) — עובד על מכשיר אמיתי, עם חידוש אחת ל-7 ימים;
  \item שלב ההפצה: חשבון מפתח בתשלום (99 דולר לשנה) מנפיק תעודות וקובצי פרופיל — רק כך אפשר לעלות לחנות ולהתקין לטווח ארוך.
\end{itemize}

אריזה: בתפריט Xcode בוחרים Product Archive (ארכיון); אחרי הבנייה פותחים את Organizer, בוחרים Distribute App App Store Connect העלאה. חשוב: לארכיון חייבים לבחור Generic iOS Device כיעד (לא סימולטור) — זה השלב שמתחילים שוכחים הכי הרבה.

\textbf{אזהרה: תעודה שאבדה = הפסקת עדכונים}\par

אם תעודת החתימה והמפתח הפרטי אבדו, לא תוכלו לשלוח עוד עדכונים לאפליקציה הישנה — המשתמשים יוכלו רק להתקין אפליקציה “חדשה לגמרי”. גבו את התעודה כראוי (ייצוא .p12 מחולץ המפתחות) ושמרו היטב על הסיסמה. למזלכם, אפל מספקת “חתימה אוטומטית” וניהול תעודות בענן (סמנו Automatically manage signing ב-Xcode) — זה חוסך את רוב המלכודות של העבודה הידנית.

\booksection{תהליך העלאה מלא ל-App Store}

רישום חשבון מפתח

הכנת חומרים לחנות

העלאת חבילת הבנייה

מילוי פרטי פרטיות ובדיקה

בדיקת Apple (1-7 ימים)

פרסום

\begin{center}\small\color{gray}איור 5: ששת שלבי העלייה ל-App Store\end{center}

\begin{itemize}
  \item הרשמה: developer.apple.com, חשבון אישי ב-99 דולר לשנה;
  \item חומרים: אייקון (1024×1024, בלי פינות מעוגלות — אפל קובעת את העיגול בעצמה), צילומי מסך (בגדלים השונים של iPhone), תיאור וקישור למדיניות פרטיות;
  \item שאלון פרטיות: מדווחים ביושר אילו נתונים נאספים — ספר חשבונות שלא מתחבר לרשת יכול לסמן "לא אוספים דבר";
  \item ביקורת: הגשה ראשונה כרוכה בהמתנה של 1–7 ימים בדרך כלל; אם נדחיתם, תקנו סעיף-סעיף לפי "הערות הביקורת" — סיבות הדחייה של אפל בדרך כלל כתובות במפורש מאוד, פשוט בצעו לפיהן;
  \item TestFlight: לפני העלייה לחנות שלחו קודם גרסת בטא לחברים (עד 10,000 איש), אספו סבב משוב ורק אז הגישו רשמית — זה נראה מקצועי בטירוף.
\end{itemize}

\booksection{חומרי חנות: הרושם הראשוני הוא ההורדות}

עמוד החנות הוא חלון הראווה שהמשתמש רואה כשהוא “מסתובב בשוק”. שלושה חומרים קובעים את הרושם הראשוני:

\begin{itemize}
  \item אייקון: 1024×1024, אל תסתפקו בצילום מסך. אין לכם כישורי עיצוב? צרו גרפיקה פשוטה עם כלי AI; האייקון עובר את המבחן רק אם הוא ניכר מבט אחד גם אחרי כיווץ על שולחן העבודה;
  \item צילומי מסך: לפחות 2 סטים (בגדלים הנפוצים של iPhone), שמציגים את תרחישי הפונקציונליות המרכזיים; לכל צילום מוסיפים משפט אחד בשפת בני אדם ("רושמים חשבון בשנייה" טוב פי עשרה מ"פתרון יעיל לניהול חשבונות");
  \item תיאור: המשפט הראשון מסביר בבירור "מה זה + למי זה מיועד"; אל תערמו "הכי מקצועי, הכי חזק" — כנות משפיעה יותר מהגזמה.
\end{itemize}

\begin{tcbpractice}{הידיים על ההגה: "מוסרים" את ספר החשבונות פעם אחת}
\booksection{רשימה קצרה: דיווח קריסות ואופטימיזציית ביצועים}

לפני שהאפליקציה עולה לחנות, שווה להשקיע חצי שעה בכל אחד משני הדברים האלה:

\booksubsection{ניטור קריסות}

זה שלא קרס אצלכם בבדיקות לא אומר שלא יקרוס בטלפון של המשתמשים. אפל מספקת דוחות קריסה חינמיים ב-Xcode Organizer (בחשבון בתשלום אפשר לראות לוגים ממכשירים אמיתיים); פתרון נוח עוד יותר הוא שילוב SDK חיצוני (כמו Firebase Crashlytics). העלייה לחנות בלי ניטור קריסות זה כמו לנהוג בעיניים עצומות.

\booksubsection{ארבעה סעיפים לבדיקת ביצועים עצמית}

\begin{enumerate}
  \item רשימות עם LazyVStack / List: אל תסדרו רשימה ארוכה ב-VStack רגיל — טעינה עצלה היא השיעור הראשון בביצועים;
  \item אל תציפו בתמונות מקור בגודל מלא: השתמשו ב-AsyncImage (תמונה מרוחקת) ובמטמון המערכת — תמונת מקור של 5MB מספיקה כדי לחנוק מכשיר ישן;
  \item אל תעשו עבודה כבדה בת׳רד הראשי: רשת, קבצים גדולים — הכל דרך async/await;
  \item מדדו עם Instruments: מנתח הביצועים המובנה ב-Xcode (Time Profiler, בדיקת דליפות זיכרון) — קודם מודדים, רק אז מייעלים; אל תסמכו על תחושת בטן.
\end{enumerate}

(1) הוסיפו 3 בדיקות יחידה ללוגיקה כמו parseAmount (מריצים עם CmdU עד לירוק); (2) תפסו בעצמכם עם נקודת עצירה קריסה שיצרתם במו ידיכם ותקנו אותה; (3) ארזו חבילה עם Product Archive והזמינו חברים לבטא דרך TestFlight; (4) אם נשאר לכם כוח — רשמו חשבון מפתח והגישו בפועל את ספר החשבונות לחנות. את שלושת הצעדים הראשונים אפשר לסיים עוד הלילה.

\end{tcbpractice}

\begin{tcbsummary}{סיכום הפרק}
\begin{itemize}
  \item שלושת כלי הניפוי: לוג print, נקודות עצירה, תצוגה מקדימה של SwiftUI ובודק התצוגות
  \item XCTest: בדיקות יחידה (CmdU) + בדיקות UI + בדיקה ידנית — התבנית כבר מוגדרת מראש
  \item חתימה: חשבון חינמי עם חתימה אוטומטית (חידוש כל 7 ימים); רק חשבון בתשלום מאפשר הפצה
  \item אריזה: Generic iOS Device + Archive + Distribute; חובה לגבות את התעודה
  \item שישה שלבי עלייה לחנות: הרשמה חומרים העלאה שאלון פרטיות ביקורת פרסום; בטא קודם עם TestFlight
\end{itemize}

\end{tcbsummary}
